\subsection{Large-Sample Confidence Intervals for a Population Mean and Proportion}
\hrulefill

\paragraph*{Confidence Intervals for Large Samples:}
If the sample size $n$ is large then the sample mean has approximately a Normal distribution regardless of the 
population distribution according to teh Cenreal Limit Theorem. However, in most practical applications 
the population variance $\sigma^2$ is unknown. It is still possible to obtain conficence intervals for the 
mean, by first estimating $\sigma^2$ by $s^2$ and then using the estimate $s$ in the conficence interval formula.
If $n$ is sufficiently large (usually $n \geq 40$), then the standardized variable
\[
    Z = \frac{\bar{X} - \mu}{\frac{S}{\sqrt{n}}} \sim \mathcal{N}(0,1)
\]
has approximately a standard normal distribution. The $n$ needs to be a little bit larger here 
because of the additional variation introduced through the estimation of $\sigma$.

\paragraph*{Definition:}
After observing a large ($n \geq 40$) random sample from a population with unknown variance $\sigma^2$ and computing the sample mean $\bar{X}$,
an approximate $100(1-\alpha)\%$ confidence interval for $\mu$ is given by
\[
    CI_\mu = (\bar{X} - z_{\frac{\alpha}{2}} \frac{s}{\sqrt{n}}, \bar{X} + z_{\frac{\alpha}{2}} \frac{s}{\sqrt{n}})
\]

\paragraph*{Example:} 
Suppose the instructor of a large calculus class tries to write an exam of medium difficulty with a mean score of 75.
Fifty students take the exam. The mean of these exam scores is 82 with a standard deviation of 18.\\
\\
a. Compute a 90\% confidence interval for the degree of difficulty of the exam.
\begin{align*}
    X &= \text{ a student's test score} \quad \bar{X} = 75 \quad n = 50 \quad \mu_X = 82 \quad \sigma_X = 18 \quad \alpha = .1\\
    CI_\mu &= (\bar{X} - z_0.05 \frac{\sigma}{\sqrt{n}}, \bar{X} + z_0.05 \frac{\sigma}{\sqrt{n}})\\
    &= (75 + 1.64 \frac{18}{\sqrt{50}}, 75 - 1.64 \frac{18}{\sqrt{50}})\\
    &= (75 + 4.175, 75 - 4.175) = (79.175, 70.825)\\
    &= (70.825, 79.175)
\end{align*}

b. In light of this information, we can conclude the exam turned out slightly easier than the professor had hoped.

\subsubsection*{Confidence Interval for a Population Proportion}
\paragraph*{Definition}
Let $p$ denote the (true, but possibly unknown) proportion of "successes" in a population. A random sample of $n$ 
individuals is selected fro the population and let $X$ denote the number of successes in the sample.
If the sample size $n$ is large enough, then $X$ has an approximately Normal distrubution. 
The most commonly used point estimate for $p$ is $\hat{p} = \frac{X}{n}$ the sample proportion of successes.

\paragraph*{Proposition:}
An approximate $100(1-\alpha)\%$ confidence interval for $p$ itself
\[
    CI_p = (\hat{p} - z_\frac{\alpha}{2} \sqrt{\frac{\hat{p}(1-\hat{p})}{n}}, \hat{p} + z_\frac{\alpha}{2} \sqrt{\frac{\hat{p}(1-\hat{p})}{n}})
\]

Here $\hat{p}$ is the sample proportion of "successes" and $1-\hat{p}$ is the same proportion of "failures."
The above confidence interval can be used in cases where the sample size $n$ is large enough such that 
$n\hat{p} \geq 10$ and $n(1-\hat{p}) \geq 10$.

\paragraph*{Example:}
Of 150 randomly chosen voters, 104 voted for canidate Z. What is a 98\% confidence interval for the 
population proportion of votes that Z gets?

\begin{align*}
    CI_Z &= (\frac{Z}{n} - z_{.01}\sqrt{\frac{\frac{Z}{n}(1-\frac{Z}{n})}{n}}, \frac{Z}{n} + z_{.01}\sqrt{\frac{\frac{Z}{n}(1-\frac{Z}{n})}{n}})\\
    &= (\frac{104}{150} + 2.33 \sqrt{\frac{\frac{104}{150}(1-\frac{104}{150})}{150}}, \frac{104}{150} - 2.33 \sqrt{\frac{\frac{104}{150}(1-\frac{104}{150})}{150}})\\
    &= (0.6933 + 0.0736, 0.6933 - 0.0736) = (0.6197, 0.7669)
\end{align*}

\paragraph*{Interval Estimation Main Ideas:}
\begin{itemize}
    \item Specify a random interval $(A, B)$ based on a random sample such that with a high probability the interval contains the parameter of interest.
    \item The probability of coverage is pre-specified and is called the confidence coefficient and is denoted by $1-\alpha$.
    \item Usually $\alpha$ is chosen to be .01, .05, or .1
    \item For a given data set, the corresponding interval ($a, b$) is called a confidence interval for the paremeter with confidence coefficient $1-\alpha$.
    \item Before the data is collected, the end points $A$ and $B$ are random variables since they depend on the random sample
\end{itemize}

\paragraph*{Construction of Confidence Intervals:}
\begin{itemize}
    \item Suppose the estimator is approximately normally distributed (usually by central limit theorem).
    \item Given a value $\alpha$, define $z_{\frac{\alpha}{2}}$ such that $P(Z > z_{\frac{\alpha}{2}}) = \frac{\alpha}{2}$, where $Z \sim \mathcal{N}(0,1)$.
    \item Define the $100(1-\alpha)\%$ confidence interval for the parameter to be.
    \item The interpretation is that for $100(1-\alpha)\%$ of the random samples of size $n$, the random interval constructed by the above procedure will contain the true value of the parameter.
    \item Note, the smaller SE, the narrower the confidence interval, hence more precise the estimate.
\end{itemize}